% =====================================================================
\section{Proposed Solution and Engineering Design}
% =====================================================================

\subsection{High-Level System Architecture}

PlateReader LR is a containerized modular monolith with four layers (Fig.~\ref{fig:architecture}).
Clients reach the service through a browser or through scripts. In the public demo, requests pass
Cloudflare Access and an outbound-only tunnel; on the local network they reach the API directly.
The application layer is one FastAPI process. Its routers handle HTTP, a shared layer validates
requests and formats errors, and four services do the work: inference, decoding, the prediction
policy and the batch worker. A data-access module owns the connection pool and all SQL. The data
layer holds the read-only model checkpoint, the PostgreSQL database and the backup files.

Uploaded images stay in memory for the length of one request. Only text results are written to the
database, and super-resolved frames are returned in the response without being stored. This keeps the
system within the dataset license.

\fig{architecture.pdf}{High-level architecture. Solid arrows are calls; dashed arrows are background
work or data copies.}{fig:architecture}

The inference service runs MF-SR-OCR (Fig.~\ref{fig:model}). A spatial transformer
\cite{jaderberg2015stn} straightens each frame, deformable convolutions \cite{zhu2019dcnv2} align
the frames to the middle one, and a branch super-resolves them by a factor of two. A residual backbone
with GroupNorm \cite{wu2018groupnorm} encodes each frame, attention fuses the five feature maps, and a
two-layer BiLSTM with a CTC head \cite{graves2006ctc} reads the sequence over $T=32$ steps and 37
classes. Training uses $\mathcal{L}=\mathcal{L}_{\mathrm{CTC}}+0.1\,\mathcal{L}_{\mathrm{SR}}$ with
EMA weights \cite{tarvainen2017mean}. At inference, a beam search of width 16 keeps only prefixes
that match one of the two plate layouts.

\fig{model.pdf}{MF-SR-OCR recognition pipeline as served by the inference service.}{fig:model}

\subsection{Core Component Design and Interfacing}

\textbf{Components.} Table~\ref{tab:components} lists the main components, their responsibility and
where they live in the code.

\begin{table}[tbp]
\centering
\caption{Core components of the backend (\code{backend/app/}).}
\label{tab:components}
\begin{tabularx}{\textwidth}{@{}P{18mm} L P{39mm}@{}}
\toprule
Component & Responsibility & Module \\
\midrule
Routers & HTTP endpoints, status codes, \code{Location} headers & \code{routers/*.py} \\
Schemas & Request and response models; the API contract & \code{schemas.py} \\
Error model & One \code{problem+json} shape with stable codes & \code{core/errors.py} \\
Inference & Loads the checkpoint once, records its SHA-256 as \code{model\_version}, pads 1--4 frames to 5, runs the forward pass under a model lock & \code{services/inference.py} \\
Decoding & Constrained beam search, top-5 suggestions, per-character confidence from CTC alignment ($\approx$9\,ms per track) & \code{services/decoding.py} \\
Prediction policy & Applies $\tau$ to set the status; maps rows to schemas & \code{services/predictions.py} \\
Batch worker & Parses a ZIP in memory, runs one job at a time in chunks of 4 tracks & \code{services/batches.py} \\
Data access & Connection pool (1--10), 15\,s statement timeout, all SQL & \code{db.py}, \code{schema.sql} \\
\bottomrule
\end{tabularx}
\end{table}

\textbf{Selective prediction.} For a track $x$ with decoded plate $y_1\ldots y_7$, the track confidence
under constrained decoding is the geometric mean of the per-character probabilities,
$c(x)=\bigl(\prod_{i=1}^{7}p(y_i\mid x)\bigr)^{1/7}$. A result is \code{accepted} when
$c(x)\ge\tau$ and \code{needs\_review} otherwise. The threshold is chosen on the validation set as
\begin{equation}
  \tau^{*}=\min\bigl\{\tau : \mathrm{Precision}(\tau)\ge 0.95\bigr\},
  \qquad \mathrm{Coverage}(\tau)=\frac{|\{x : c(x)\ge\tau\}|}{N},
\end{equation}
which keeps as many tracks automatic as the precision target allows. The current value
$\tau=0.90$ is provisional; with it, 887 of 999 validation tracks are accepted and 112 go to review.

\textbf{Interfaces.} Clients talk to the service through REST over HTTP under \code{/api/v1}: multipart
uploads in, JSON out, CSV for export, and one error format (RFC~9457). The 14 operations group into
three resources (Fig.~\ref{fig:apimap}). The full OpenAPI 3.1 description and one page per endpoint
are in the repository (\code{api-docs/}); the figures below show the contract and its behavior
instead of reproducing that reference.

\fig{apimap.pdf}{REST resource map of \code{/api/v1} with success status codes.}{fig:apimap}

Fig.~\ref{fig:seqsingle} shows a single-track request. Validation runs before any work, the forward
pass runs in a thread pool under the model lock, and a database connection is borrowed only to insert
the result. Fig.~\ref{fig:seqbatch} shows a batch: the API answers \code{202 Accepted} at once, a
single worker processes the job in chunks of four tracks, and the client polls for progress.
Fig.~\ref{fig:states} gives the lifecycle of a batch job and of a prediction.

\fig{seq_single.pdf}{Sequence of a single-track request (\code{POST /predictions}).}{fig:seqsingle}

\fig{seq_batch.pdf}{Sequence of an asynchronous batch job (\code{POST /batches}). The worker loop
and the client polling run in parallel.}{fig:seqbatch}

\fig{states.pdf}{State machines of a batch job (a) and of a prediction (b).}{fig:states}

\textbf{Data model.} The schema has three tables (Fig.~\ref{fig:erd}). The \code{reviews} table is
append-only: a review is a new row, so two reviewers never overwrite each other, the model's original
answer stays in \code{predictions.plate}, and model accuracy can still be computed after review.
Batch statistics are counted at read time, so no counter row is shared between writers.

\fig{erd.pdf}{Entity--relationship diagram (PostgreSQL 16). Filled diamonds mark the owning side of
an \code{ON DELETE CASCADE} relation.}{fig:erd}

\textbf{Concurrent reads and writes.} Table~\ref{tab:concurrency} lists the six problems the mixed
workload can cause and how each one is handled. Fig.~\ref{fig:concurrency} shows the resulting timing:
a single request waits for at most one chunk, and readers are not blocked by writers.

\begin{table}[tbp]
\centering
\caption{Concurrency problems and the measures implemented in Week~2.}
\label{tab:concurrency}
\begin{tabularx}{\textwidth}{@{}P{28mm} L@{}}
\toprule
Problem & Measure \\
\midrule
Database-wide write lock & PostgreSQL with MVCC instead of SQLite \\
Connection exhaustion & Pool of 1--10 connections; a request waits up to 10\,s and then receives \code{503 database\_busy} with \code{Retry-After}; no connection is held during inference \\
Model starvation & One batch worker, chunks of 4 tracks, lock released after each chunk; on CPU, 32 tracks per chunk were not faster than 1 ($\approx$0.66\,s per track) \\
Lost updates on review & Append-only \code{reviews}; the newest row is the final plate \\
Hot counter row & Batch counts computed with \code{COUNT(*) FILTER} at read time, with indexes \\
Jobs stuck after a restart & On startup, jobs left \code{queued} or \code{running} are marked \code{failed} with a reason \\
\bottomrule
\end{tabularx}
\end{table}

\fig{concurrency.pdf}{Timing of a batch job, a single request (A), database writes and UI polling.
The chunk length follows the measured 0.66\,s per track on CPU.}{fig:concurrency}

\subsection{Software Engineering Methodology and DevOps Pipeline}

\textbf{Methodology.} Work follows Kanban \cite{anderson2010kanban} with a weekly cadence
(Fig.~\ref{fig:kanban}). GitHub Issues hold the backlog, each item is sized to one week or less, at
most two items are in progress, and each item goes through one branch and one pull request. The
weekly report to the supervisor serves as the review and re-planning meeting. Kanban suits a single
developer whose week mixes long GPU runs with short coding tasks, where fixed sprints would add
ceremony without a team to coordinate.

\fig{kanban.pdf}{Kanban board layout and working rules. Cards show example items around Week~3.}{fig:kanban}

\textbf{Model validation protocol.} Model results follow fixed rules (Fig.~\ref{fig:protocol}). Every
configuration is trained with three seeds; data splits are stored with SHA-256 hashes; the threshold
is chosen on validation data; and the sealed test set is used exactly once, after the model, its
configuration and $\tau$ are frozen. A parity check against the offline model is part of every
release.

\fig{protocol.pdf}{Data split (to scale) and model validation protocol.}{fig:protocol}

\textbf{CI/CD.} GitHub Actions \cite{github_actions} will run the pipeline in Fig.~\ref{fig:cicd}.
Every pull request runs static checks, unit tests, integration tests against a PostgreSQL service
container, and security scans. A release tag adds a gate on the lab machine (parity and latency),
builds and pushes the images, deploys them with Compose and runs a smoke test, with rollback to the
previous tag if the smoke test fails. The workflows are scheduled for Weeks~3--4; until then the same
checks are run by hand before each merge.

\fig{cicd.pdf}{CI/CD pipeline on GitHub Actions.}{fig:cicd}

\textbf{Deployment.} Docker Compose runs the stack on one Linux host without a GPU
(Fig.~\ref{fig:deploy}). The database port is bound to \code{127.0.0.1}, the API container runs as a
non-root user with a 3\,GB memory limit, and the checkpoint is mounted read-only, so changing the
model is a configuration change. The tunnel and backup services sit in Compose profiles and start only
when needed. The \code{web} service is added in Week~8.

\fig{deploy.pdf}{Deployment view (Docker Compose on one CPU host).}{fig:deploy}
